\anexo{Transcripción de la entrevista — Familiar cuidador a distancia}

\noindent\textbf{Perfil:} Hijo/a o familiar que cuida a un adulto mayor que vive en otra casa.\\
\noindent\textbf{Modalidad:} Semiestructurada $|$ Duración: 30 min aproximadamente.

\bigskip
\textbf{Bloque 1 — Contexto}

\begin{description}[leftmargin=0cm, labelsep=0.5cm]

    \item[\textbf{Entrevistador:}] ¿Podés contarme un poco sobre tu mamá? ¿Cómo es su día a día, con quién vive, cómo está de salud?

    \item[\textbf{Entrevistada:}] Vive sola. Físicamente está muy bien de salud, pero su memoria está un poco deteriorada. Es una persona muy activa, aunque con frecuencia suele olvidarse de las cosas. En su día a día limpia la casa y le dedica gran parte del tiempo a las plantas. No le gusta salir a hacer los mandados y le cuesta socializarse. Suele perder la noción del tiempo.

    \item[\textbf{Entrevistador:}] ¿Con qué frecuencia la ves en persona o hablás con ella? ¿Cómo se comunican?

    \item[\textbf{Entrevistada:}] Hablamos todos los días por llamada telefónica, ya que no maneja WhatsApp ni otro medio de comunicación. Día por medio paso por su casa a visitarla, aunque sean unos minutos, para ver cómo está y si necesita algo.

\end{description}

\bigskip
\textbf{Bloque 2 — El problema}

\begin{description}[leftmargin=0cm, labelsep=0.5cm]

    \item[\textbf{Entrevistador:}] ¿Hay momentos del día en los que te preguntás cómo estará o qué estará haciendo?

    \item[\textbf{Entrevistada:}] Sí, casi siempre a la mañana me pregunto si ya se levantó. Al mediodía pienso si está cocinando, si no se le quema la comida, si está comiendo. También suelo pensar en ella cuando vuelvo del trabajo, tanto al mediodía como a la tarde-noche.

    \item[\textbf{Entrevistador:}] ¿Alguna vez pasó algo que te haya generado preocupación o que te haya tomado desprevenida?

    \item[\textbf{Entrevistada:}] Sí. Una vez le dejé varios palitos Bom Bon helado en el freezer para que tuviera para comer de postre de vez en cuando, y ella se comió toda la caja en dos días. Eso generó que al otro día sufriera un ataque de hígado. Estuvo con vómitos sin contarle a nadie hasta que realmente no aguantó más. Cuando nos avisó, tuvimos que llevarla al hospital, donde quedó internada unas horas.

    \item[\textbf{Entrevistador:}] ¿Qué es lo que más te preocupa cuando no estás presente?

    \item[\textbf{Entrevistada:}] Principalmente que cuando cocine no prenda fuego la casa, que no le abra la puerta a ningún extraño y que siempre cierre con llave. Si bien está bien físicamente, al ser muy activa a veces quiere hacer cosas que no son para su edad y pueden llevarla a sufrir una caída o provocarle algún dolor. Por ejemplo, se sube arriba de la silla para bajar cosas de arriba de la alacena.

\end{description}

\bigskip
\textbf{Bloque 3 — Soluciones actuales}

\begin{description}[leftmargin=0cm, labelsep=0.5cm]

    \item[\textbf{Entrevistador:}] ¿Hoy en día hacés algo para estar más tranquila? ¿Usás alguna app, algún dispositivo?

    \item[\textbf{Entrevistada:}] Tiene celular. Con mis hermanos nos turnamos para ir a visitarla, pero la realidad es que no usamos ningún dispositivo de control. Sí pensamos en instalarle cámaras.

    \item[\textbf{Entrevistador:}] ¿Qué funciona bien y qué no funciona?

    \item[\textbf{Entrevistada:}] Lo que no funciona es que ella no logra hacer videollamadas, no siempre atiende el celular y tampoco puedo enviarle mensajes porque no maneja WhatsApp. Lo que funciona es ir a visitarla: ahí nos ponemos al día de qué estuvo haciendo, qué va a hacer, qué comida tiene, qué le falta y cómo está.

\end{description}

\bigskip
\textbf{Bloque 4 — Validación de la propuesta}

\begin{description}[leftmargin=0cm, labelsep=0.5cm]

    \item[\textbf{Entrevistador:}] Si existiera un sistema que detecta automáticamente si tu mamá se cayó o lleva mucho tiempo sin moverse, y te avisa por el celular, ¿qué te parecería?

    \item[\textbf{Entrevistada:}] Sí, excelente. Es realmente lo que necesitamos. Un sistema que nos haga sentir más cerca de ella.

    \item[\textbf{Entrevistador:}] El sistema usa cámaras en el hogar pero todo el video se procesa ahí mismo, no sale de la casa. ¿Eso te genera alguna duda?

    \item[\textbf{Entrevistada:}] Me parece bien, siempre y cuando no se grabe al adulto mayor en lugares de mayor privacidad, como el baño. Fuera de eso, me parece una gran idea.

\end{description}